\documentclass[11pt]{article}
\usepackage[margin=0.8in]{geometry}
\usepackage[T1]{fontenc}
\usepackage{lmodern}
\usepackage{booktabs}
\usepackage{longtable}
\usepackage{microtype}
\usepackage{hyperref}
\usepackage{xcolor}
\hypersetup{colorlinks=true,linkcolor=blue!50!black,urlcolor=blue!50!black}
\setlength{\parindent}{0pt}
\setlength{\parskip}{5pt}

\title{\textbf{Measuring the Local Weight of Coding-Agent Harnesses}\\
\large A controlled CPU, memory, and task-deployment comparison\\
\normalsize with an OpenRouter Grok 4.6 supplement}
\author{mmdmcy\\Katteke}
\date{20 August 2026}

\begin{document}
\maketitle

\begin{abstract}
Coding-agent harnesses are often compared by model quality or task speed, but a
developer running many agents also pays for the local process tree: resident
memory, CPU time, child processes, and startup work. This thesis reports a
controlled follow-up study of OpenCode, Codex CLI, Cursor Agent CLI, Pi, xAI
Grok Build, and Vercel fx. The primary endpoints were peak process-tree resident
memory and total user plus system CPU seconds. Twenty fresh version/help trials
per harness measured cold startup. Three repeated light and heavyweight coding
trials measured deployment behavior, with solo and delegation-enabled arms where
the CLI exposed a comparable mechanism.

fx was the lightest cold-start binary at a median 1.47 MB peak RSS. In the
matched GPT-5.6 Luna-max solo task subset, Pi had the lowest median peak tree
RSS on both light and heavy tasks, Codex used the least CPU on the heavy task,
and Cursor finished that task fastest. No harness was best on every metric.
These are narrow deployment measurements: task results include model/provider
behavior and tool topology, while fx and Grok use native models rather than the
matched Luna-max family. Grok task repetition was censored by its free usage
quota, and fx timed out on all three heavyweight trials. A separate Grok Build
supplement routed Grok 4.6 through OpenRouter: all 12 runs completed, with solo
heavy medians of 220.4 MB peak RSS, 1.41 CPU seconds, 380.7 wall seconds, and
23/24 hidden checks.
\end{abstract}

\section{Question and Estimands}

The motivating question was not simply which CLI feels fastest. It was which
installed harness places the smallest local burden on a four-core Linux machine,
and whether that answer changes for ordinary maintenance work, heavyweight
multi-package work, and delegation. The study therefore uses two estimands.

The first estimand is fresh-process startup footprint. Each executable ran
\texttt{--version} and \texttt{--help} twenty times in a seeded sequential
order. No model request was made. The second is deployment footprint: the full
local process tree created by an agent while it attempted a fixed coding task.
The second estimand deliberately includes model response handling and tool use;
it answers a practical deployment question rather than isolating binary code
size from all other work.

\section{Systems and Controls}

The matched task arms explicitly requested GPT-5.6 Luna max wherever available:
OpenCode used \texttt{openai/gpt-5.6-luna} with variant \texttt{max}; Codex used
\texttt{gpt-5.6-luna} with \texttt{model\_reasoning\_effort=max}; Cursor used
\texttt{gpt-5.6-luna-max}; and Pi used \texttt{openai-codex/gpt-5.6-luna} with
\texttt{--thinking max}. Grok Build used native Grok 4.6. fx used its native
default \texttt{zai/glm-5.2}. The latter two are explicitly separate model
exceptions rather than hidden claims of model equivalence.

The model distinction matters even though all inference is remote. Different
models can choose different numbers of turns, tool calls, subprocesses, and
delegation patterns; those choices alter the measured local process tree. The
Luna-max subset is therefore the closest harness comparison. The fx, native
Grok, and OpenRouter Grok arms are valid deployment-system observations but not
pure causal estimates of harness implementation overhead.

The host was Linux x86-64 with four logical CPUs and approximately 15.5 GiB of
RAM. Runs were sequential. Each task received a fresh fixture copy. The light
fixture was a small retry-with-backoff maintenance issue with eight hidden
checks. The heavy fixture contained three independent implementation, security,
and comprehension packages with 24 hidden checks. All agent work was scored
after the resource measurement window.

Solo arms disabled delegation where possible. Delegated arms enabled the
harness-specific mechanism: bounded OpenCode workers, Codex multi-agent mode,
Cursor task tools, and a Pi Luna-worker extension. fx exposed no comparable
subagent control in this version, so it has a solo arm only.

\section{Instrumentation}

The runner samples Linux \texttt{/proc} every 20 ms and follows descendants of
the launched CLI process. It records summed tree RSS, maximum individual RSS,
CPU ticks, process count, thread count, context switches, and process I/O. Very
short native binaries can exit before a sampling interval; startup commands
therefore also run under GNU \texttt{time -v}, with child rusage supplying the
fallback CPU and maximum-RSS values. This fallback is applied consistently to
the startup phase.

The primary summaries are medians and interquartile ranges. Task cells contain
only three trials, so no significance test or universal rank is claimed. A
provider quota failure is classified as provider-blocked. A task timeout remains
an observed timeout and is not silently deleted.

\section{Cold Startup Results}

\begin{center}
\begin{tabular}{lrrr}
\toprule
Harness & Peak RSS & CPU seconds & Wall seconds \\
\midrule
fx & 1.47 MB & $<0.01$ & 0.21 \\
Grok Build & 27.0 MB & 0.02 & 0.22 \\
Codex CLI & 50.5 MB & 0.09 & 0.24 \\
Pi & 157.9 MB & 2.51 & 2.36 \\
OpenCode & 187.5 MB & 1.85 & 1.67 \\
Cursor Agent & 206.4 MB & 1.76 & 1.70 \\
\bottomrule
\end{tabular}
\end{center}

The result supports the strongest narrow conclusion of this study: fx is the
lightest fresh native CLI in this environment. It does not establish that an
active fx coding session remains the smallest after model context, tools, and
long-running work are included.

\section{Matched Luna-Max Task Results}

\begin{center}
\begin{tabular}{lrrrr}
\toprule
\multicolumn{5}{c}{Light maintenance, solo, three trials} \\
Harness & Peak RSS & CPU s & Wall s & Score \\
\midrule
Pi & 240.6 MB & 3.70 & 52.6 & 8/8 \\
Cursor & 363.4 MB & 6.23 & 59.1 & 8/8 \\
Codex CLI & 441.1 MB & 4.79 & 72.7 & 8/8 \\
OpenCode & 742.8 MB & 19.57 & 81.0 & 8/8 \\
\bottomrule
\end{tabular}
\end{center}

\begin{center}
\begin{tabular}{lrrrr}
\toprule
\multicolumn{5}{c}{Heavy independent packages, solo, three trials} \\
Harness & Peak RSS & CPU s & Wall s & Median score \\
\midrule
Pi & 535.8 MB & 18.86 & 1047.5 & 22/24 \\
Cursor & 565.0 MB & 18.53 & 458.3 & 22/24 \\
Codex CLI & 608.9 MB & 12.92 & 851.7 & 22/24 \\
OpenCode & 899.6 MB & 49.52 & 836.5 & 22/24 \\
\bottomrule
\end{tabular}
\end{center}

The endpoints disagree by design. Pi is the lightest by peak tree RSS in both
matched task types. Codex is the CPU-minimal heavy-task system. Cursor is the
wall-time-minimal heavy-task system. The data do not support a single scalar
claim that one harness is simply ``lightest.''

\section{Delegation}

Delegation changed topology and usually increased heavy-task CPU. The median
delegated-to-solo CPU ratios were 1.36 for OpenCode, 2.41 for Codex, 2.22 for
Cursor, and 2.31 for Pi. Heavy-task wall ratios were 1.14, 1.17, 1.30, and
1.07 respectively. RSS ratios were not monotonic because child overlap and
parent context release determine the peak.

This is not evidence that delegation is intrinsically wasteful. Delegation arms
were asked to do more parallel work, and several scored 22--23/24 rather than a
perfect 24/24. The correct interpretation is that the measured process tree
price of the delegated deployment was generally higher in CPU, while the speed
and memory consequences depended on the harness.

\section{Native Exceptions}

fx completed all three light tasks at 8/8. Its three heavy trials reached the
1800-second task limit with scores 2/24, 14/24, and 2/24. Its low local resource
envelope is therefore not a heavy-task efficiency win. Grok Build completed one
light solo task and one delegated heavy task before the free usage quota was
exhausted; eleven later task attempts were provider-blocked. Grok is reported
as a native-harness observation, not ranked against the repeated Luna arms.

\section{OpenRouter Grok Supplement}

To separate the native Grok Build quota from the harness, a post-hoc supplement
configured Grok Build 1.0.5 with the OpenAI-compatible OpenRouter endpoint and
model \texttt{x-ai/grok-4.6}. The API key was supplied only through an
environment variable. The isolated config, aggregate results, and manifest do
not contain the key. The task limit was increased from 1800 to 3600 seconds and
the turn limit from 36 to 64 before these supplement outcomes were collected.

\begin{center}
\begin{tabular}{llrrrr}
\toprule
Architecture & Task & Peak RSS & CPU s & Wall s & Score \\
\midrule
Solo & Light & 213.9 MB & 0.54 & 27.0 & 8/8 \\
Delegated & Light & 163.6 MB & 0.60 & 29.1 & 8/8 \\
Solo & Heavy & 220.4 MB & 1.41 & 380.7 & 23/24 \\
Delegated & Heavy & 307.7 MB & 6.92 & 692.6 & 23/24 \\
\bottomrule
\end{tabular}
\end{center}

All 12 planned supplement runs completed. This supports a practical conclusion:
the earlier repeated-run failure was an xAI free-service quota limitation, not
evidence that the open-source Grok Build harness could not perform the tasks.
It does not make the supplement directly interchangeable with the Luna-max arms.
The provider route, model, timeout, and turn budget differ, so the correct label
is \emph{Grok Build with OpenRouter Grok 4.6}.

Delegation was not beneficial on this heavy fixture in the supplement. Relative
to solo medians, it used about 40\% more peak RSS, 391\% more CPU, and 82\% more
wall time while retaining the same 23/24 median score. With three trials per
cell, this is descriptive rather than a general conclusion about subagents.

\section{Validity and Limitations}

The design is valid for the narrow operational question of observed local
resource envelopes under this host, these versions, these tasks, and these
provider configurations. It is not a universal benchmark of harness quality or
model efficiency. Limitations include one machine and service window, three
task repetitions, cloud-provider behavior, a synthetic separable heavy fixture,
no GPU or energy measurement, no persistent IDE-session measurement, unequal
delegation implementations, and model exceptions for fx and Grok. The fixed
process-tree instrumentation also cannot observe provider-side compute. The
20-ms sampler can miss very short descendant processes; GNU \texttt{time}
provides a startup fallback but not full descendant accounting for task runs.
Task CPU should therefore be read as measured local CPU, potentially a lower
bound for very short-lived children.

The task pair is inspired by common coding-agent evaluation shapes but is not
SWE-bench, Terminal-Bench, or another official benchmark. The report therefore
uses descriptive medians and explicit availability dispositions instead of
pretending to supply a general industry leaderboard.

\section{Conclusion}

If the question is cold-start footprint, fx is the clear local minimum in this
sample. If the question is a repeated GPT-5.6 Luna-max coding deployment, Pi
used the least peak process-tree memory, Codex used the least CPU on the heavy
task, and Cursor completed that heavy task fastest. OpenCode was the heaviest
local tree in this particular configuration. The appropriate engineering choice
depends on whether memory, CPU, latency, task completion, or delegation
capacity is the binding constraint.

The OpenRouter supplement adds a strong operational result: Grok Build with
OpenRouter Grok 4.6 completed every task cell with a small local process tree,
and its solo heavy median was 220.4 MB RSS and 380.7 seconds at 23/24. It is the
best observed heavy deployment in this dataset on those two endpoints, but not
the winner of a model-controlled harness experiment. A future causal harness
comparison should run the same OpenRouter model, endpoint, prompt, tool policy,
timeout, and effort setting through every harness that supports it.

\end{document}
